\subsection{Prediction and averaging on marked cells}
\label{pred:section}

Fix a shift at a vertex of the analytic parameter simplex, with imaginary
parameters of Euclidean norm at most \(L^2\). Retain the block notation,
path top \(R\), and meshes from the preceding construction. Thus \(J\) is the
path interval at level \(i_j\), its block children \(c\) have width
\(\rho_{j+1}\), and
\[
 G\rho_{j+1}\le\rho_j\le B\rho_{j+1},\qquad
 H_j=\frac{w}{\rho_{j+1}\sqrt G},\qquad
 H_{j-1}=\frac{w}{\rho_j\sqrt G}.
\]
The first inequalities include the initial, possibly shorter, block.
The path child is denoted by \(a\) when it exists. At the last path
block there is no path child, even when some other tops continue; in that
case the relative path drop is \(u_j=1\).

\begin{proposition}[Averaging factors on marked cells]
\label{pred:marked-comparison}
There is an absolute \(c_{\rm av}>0\) with the following property, for
all sufficiently large \(L\). On a fixed line of slope \(s_R\), each
marked \(H_j\)-cell \(I\) meeting \(Z\) admits a nonnegative function
\(\Gamma_j\) which is constant on the \(H_{j-1}\)-subcells of \(I\),
satisfies \( |I|^{-1}\int_I\Gamma_j\le1\), and obeys
\begin{equation}
 \beta_J^0(z)\le
       \exp\bigl(G^{-c_{\rm av}}u_j(z)\bigr)\Gamma_j(z)
       \qquad (z\in I).
 \label{eq:src-24}
\end{equation}
Set \(\Gamma_j=1\) on unmarked cells and on cells not meeting \(Z\).
The functions can be chosen measurably on each fixed line. No choice
jointly measurable in the family of parallel lines is needed.
\end{proposition}

We first construct a normalized prediction \(\widetilde\beta\) whose
average is at most one for suitable probability measures. We compare it
with \(\beta_J^0\), then blend it with one to obtain \(\mathcal B\), whose
comparison error contains the relative path drop. Finally we pass to
uniform cell measure, take the finer-cell supremum \(\mathcal B^\sharp\),
and normalize it to obtain \(\Gamma_j\). The losses in these last
operations also contain the relative path drop.

\subsection{Snapshot data and denominator shrinkage}

Fix a marked cell \(I\) meeting \(Z\), and choose one snapshot
\(z_0\in I\cap Z\). Parameterize the line by horizontal displacement
\(b\) from \(z_0\), and also write \(I\) for the translated parameter
interval. All unprimed masses and amplitudes below are evaluated at
\(z_0\); a prime denotes evaluation at \(z_0+(b,s_Rb)\).
Let \(n=n_J(z_0)\) and \(\upsilon=u_j(z_0)\).
If the path child exists, then \(\upsilon=1-n_a/n\); otherwise
\(\upsilon=1\). The positive-mass motion bounds give
\begin{equation}
 \tfrac12\upsilon\le u_j(z)\le2\upsilon\qquad(z\in I)
 \label{pred:drop-comparability}
\end{equation}
for large \(L\). Marking ensures \(\upsilon>0\).
Only nonempty children are used. Their masses are positive at every
point, since the sets of eligible tops are fixed and their weights are
strictly positive.

Use the pointwise normalization \eqref{eq:src-25}, with \(l_c\) and
\(X_c\) evaluated at the snapshot and the primed quantities evaluated
at the displaced point. That identity expresses \(\beta_J^0\) in either
set of data. The next construction shrinks snapshot vectors to obtain
one-sided denominator comparisons for those children. When
\(\upsilon<1/2\), the path vector will instead remain unchanged so that
its numerator and denominator variations can be compared together.
All vector norms in this section are in \(\mathcal H\).
The weights satisfy \(\sum_c l_c\le1\), and, when there is a path
child, \(\sum_{c\ne a}l_c\le\upsilon\). When there is none,
an off-path sum means the sum over all children.

Choose a fixed exponent \(C_*\), large enough for the estimates below,
and put
\begin{equation}
 K_*=L^{C_*}(1+N^\mu).
 \label{pred:amplitude-scale}
\end{equation}
Equations~\eqref{eq:src-12} and \eqref{eq:src-17} imply, at the
snapshot,
\begin{equation}
 \norm{X_c}\le K_*,\qquad n^\mu\le K_*,\qquad
 \sum_{c\ne a}l_c^{1-\mu}\norm{X_c}\le K_*\upsilon.
 \label{pred:snapshot-size}
\end{equation}
Enlarging the fixed exponent absorbs absolute constants here and below.
For the last estimate use
\[
 l_c^{1-\mu}X_c=n^{\mu-1}A_c
              =n^\mu l_c(A_c/n_c).
\]
The sum over all children in this estimate is at most \(K_*\).

Let \(\gamma>0\) be a sufficiently small fixed exponent below the
motion exponent in \Cref{hier:amplitude-motion}, and set
\(\varepsilon=G^{-\gamma}\). This choice is independent of the depth,
dimension, and parameters \(L,p,D\). Decreasing \(\gamma\), if
necessary, absorbs the fixed powers of \(L\) in the following
consequences of the motion estimates:
\begin{align}
 \left\|l_c'^{1/p}X_c'
       -\mathcal D_c(b)l_c^{1/p}X_c\right\|
       &\le\varepsilon l_c^{1/p}n_c^\mu,
       \label{pred:scaled-motion}\\
 \left\|n'^{\mu-1}\lambda_c'A_c'
       -n^{\mu-1}\lambda_c\mathcal D_c(b)A_c\right\|
       &\le\varepsilon K_*l_c.
       \label{pred:numerator-motion}
\end{align}
Both estimates hold throughout \(I\), without a good-point hypothesis
at the displaced point. For the scalar rescaling in the first one,
write \(l_c^{1/p}X_c=n^{-1/p}n_c^{\alpha/p-1}A_c\).
The exponents \(-1/p\) and \(\alpha/p-1\) are absolutely bounded.
Their factors have the stated relative mass motion; the error in
\(A_c\) is bounded by \eqref{eq:src-21}, and its snapshot size is at
most \(L^C n_c\). This proves \eqref{pred:scaled-motion}.
For the second estimate use \(n^{\mu-1}\), the relative complex motion
of \(\lambda_c\), \(\abs{\lambda_c}\le1\), and
\(n^{\mu-1}n_c=n^\mu l_c\le K_*l_c\).

Shrink \(A_c\) radially toward zero by a distance at most
\(2\varepsilon n_c\), choosing
\[
 \widetilde A_c=
 \begin{cases}
   (1-2\varepsilon n_c/\norm{A_c})_+ A_c,& A_c\ne0,\\
   0,& A_c=0.
 \end{cases}
\]
If no eligible top in \(c\) is significant at \(z_0\), set
\(\widetilde A_c=0\). This fits within the same error:
\Cref{hier:amplitudes} and the mask suppression give
\[
 \norm{A_c}\le L^C e^{-L\log L}n_c\le\varepsilon n_c
\]
for such a child. There is one exception: when \(\upsilon<1/2\),
leave the path vector unchanged, setting \(\widetilde A_a=A_a\).
Define \(\widetilde X_c=n_c^{\mu-1}\widetilde A_c\).
For every child to which shrinkage was applied,
\begin{equation}
 l_c'\norm{X_c'}^p\ge l_c\norm{\widetilde X_c}^p
             \qquad(b\in I).
 \label{pred:denominator-lower-bound}
\end{equation}
Indeed radial shrinkage reduces the norm of \(l_c^{1/p}X_c\) by
the smaller of its whole norm and \(2\varepsilon l_c^{1/p}n_c^\mu\),
whereas \eqref{pred:scaled-motion} permits a reduction of at most
\(\varepsilon l_c^{1/p}n_c^\mu\) at a displaced point.

The bounds in \eqref{pred:snapshot-size} continue to hold with
\(\widetilde X_c\). Including the shrinkage error in
\eqref{pred:numerator-motion} gives
\begin{equation}
 \left\|l_c'^{1-\mu}\lambda_c'X_c'
       -l_c^{1-\mu}\lambda_c\mathcal D_c(b)\widetilde X_c\right\|
       \le C\varepsilon K_* l_c.
 \label{pred:shrunk-numerator-motion}
\end{equation}

\subsection{Exact averaging of the prediction}

On the whole parameter line, define
\begin{equation}
 F(b)=\sum_c l_c^{1-\mu}\lambda_c
                   \mathcal D_c(b)\widetilde X_c,\qquad
 \widetilde\beta(b)=
       \frac{1+\norm{F(b)}^p}
            {1+\sum_c l_c\norm{\widetilde X_c}^p}.
 \label{pred:model}
\end{equation}
The masses, scalar weights, and vectors are fixed snapshot data.
Only the diagonal unitaries
\(\mathcal D_c(b)_{\xi\xi}=e^{i(s_R-t_c)\eta_\xi b}\) vary.

\begin{lemma}[Fourier-supported averaging]
\label{pred:fourier-average}
There is a fixed sufficiently small \(\epsilon'>0\) such that, with
\begin{equation}
 Q=G^{3/4}\rho_{j+1},\qquad \delta=\epsilon' Q/w,
 \label{pred:separation-bandwidth}
\end{equation}
every positive probability \(\nu\) satisfying
\(\supp\widehat\nu\subset(-\delta,\delta)\) obeys
\(\int\widetilde\beta\,\dd\nu\le1\).
\end{lemma}

\begin{proof}
Use only children for which \(\widetilde X_c\ne0\); an empty
collection gives the assertion immediately. Apply the clock selection
from \eqref{pred:retention-probability} to this finite subfamily at the
snapshot, with \(Q\) as in \eqref{pred:separation-bandwidth}. The clock
rates are fixed as \(b\) varies. Each retained set has pairwise interval
distance at least \(Q\), and the retention probabilities satisfy
\(\pi_c\ge n_c/m_c\).

Every used child other than the possible unshrunk path child contains
a significant eligible top at \(z_0\). From each retained child other
than that possible exception, select one such top. Their slopes are
\(Q\)-separated, their
precisions satisfy
\[
 d_T\le\rho_{j+1}/B^2\le Q/B^2,
\]
and their diameter is at most
\(\rho_j\le B\rho_{j+1}\le BQ\).
The simultaneous count in \Cref{count:simultaneous}, with \(D\) as
chosen in \Cref{count:choice-D}, bounds this list by \(D/100\).
Adding the possible retained path child leaves at most \(D\) members.
Clock exclusion separates that extra child too; it does not require
a significant representative for that one child.

In each Hilbert component, the retained carrier frequencies are
separated by at least \(c Q/w\), with \(c>0\) absolute, since
\(\eta_\xi\asymp w^{-1}\). Choose \(\epsilon'\le\epsilon_0c\).
For a retained set \(\mathcal S\), let
\[
 F_{\mathcal S}(b)=
     \sum_{c\in\mathcal S}\pi_c^{-1}
       l_c^{1-\mu}\lambda_c\mathcal D_c(b)\widetilde X_c.
\]
Expectation over the clocks gives
\(\mathbb E_{\mathcal S}F_{\mathcal S}=F\); no independence of
retention indicators is asserted or needed. Put \(r=r_*\).
Minkowski in that expectation, \eqref{eq:src-13}, and Jensen for
the concave function \(t^{1/r}\) give
\begin{align}
 \norm F_{L^p(\nu;\mathcal H)}
 &\le \mathbb E_{\mathcal S}
       \left(\sum_{c\in\mathcal S}
          \pi_c^{-r}l_c\abs{\lambda_c}^r
                           \norm{\widetilde X_c}^r\right)^{1/r}
       \notag\\
 &\le\left(\sum_c\pi_c^{1-r}l_c\abs{\lambda_c}^r
                           \norm{\widetilde X_c}^r\right)^{1/r}
       \notag\\
 &\le\left(\sum_c l_c\norm{\widetilde X_c}^r\right)^{1/r}
 \le\left(\sum_c l_c\norm{\widetilde X_c}^p\right)^{1/p}.
 \label{pred:weighted-phase-bound}
\end{align}
The power of \(l_c\) uses \(r(1-\mu)=1\).
The third line uses the retention cancellation
\eqref{hier:selection-cancellation} for the present child subfamily.
The last inequality in \eqref{pred:weighted-phase-bound} is H\"older
with weights of total mass at most one and \(r\le p\); the missing
mass can be placed at a zero vector. Raise to power \(p\), add one,
and divide by the constant denominator in \eqref{pred:model}.
\end{proof}

\Cref{fig:phase-scales} summarizes the horizontal scales used in the
comparison and discretization below.
\begin{figure}[tbp]
\centering
\begin{tikzpicture}[x=1cm,y=1cm,font=\small,>=Stealth]
 \draw[->,thick] (0,0)--(11.7,0);
 \foreach \xx in {.7,3.8,6.75,10.5}
   \draw[thick] (\xx,-.12)--(\xx,.12);
 \node[above,align=center] at (.7,.2)
   {\(H_{j-1}\)\\[-1pt]\(\displaystyle\frac{w}{\rho_j\sqrt G}\)};
 \node[above,align=center] at (3.8,.2)
   {phase scale\\[-1pt]\(\displaystyle\frac{w}{Q}\)};
 \node[above,align=center] at (6.75,.2)
   {\(H_j\)\\[-1pt]\(\displaystyle\frac{w}{\rho_{j+1}\sqrt G}\)};
 \node[above,align=center] at (10.5,.2)
   {\(l_*\)\\[-1pt]\(\displaystyle\frac{B^2w}{\rho_{j+1}}\)};
 \draw[<->,blue!65!black] (3.8,-.65)--(6.75,-.65);
 \node[below,blue!65!black] at (5.275,-.68)
   {\(\displaystyle\frac{QH_j}{w}=G^{1/4}\)};
 \draw[<->,green!40!black] (6.75,-1.35)--(10.5,-1.35);
 \node[below,green!40!black] at (8.625,-1.38)
   {\(\displaystyle\frac{l_*}{H_j}=B^2\sqrt G\)};
\end{tikzpicture}
\caption{The horizontal-length axis (schematic) for a marked block.
Here \(Q=G^{3/4}\rho_{j+1}\) is the retained angular separation and
\(G\rho_{j+1}\le\rho_j\le B\rho_{j+1}\), including the initial block.
The quantity \(l_*=B^2w/\rho_{j+1}\) is the lower bound for predictable
packet lengths in \eqref{hier:predictable-length}.
On an interval of length \(H_j\), the separated relative phases
oscillate many times while the packet envelopes move little.
The preceding mesh also satisfies
\(H_{j-1}\rho_j/w=G^{-1/2}\), which controls the step supremums.
Absolute frequency constants are suppressed.}
\label{fig:phase-scales}
\end{figure}


\subsection{Comparison to the actual factor}

On the underlying real Hilbert space put
\(\mathcal L(Y)=\log(1+\norm Y^p)\). For \(p\ge4\),
\begin{equation}
 \norm{\nabla\mathcal L(Y)}\le p,\qquad
 \norm{D^2\mathcal L(Y)}\le C p^2.
 \label{pred:log-derivatives}
\end{equation}
Indeed its radial first derivative is \(p r^{p-1}/(1+r^p)\),
its radial second derivative is
\[
 \frac{p(p-1)r^{p-2}}{1+r^p}
       -\frac{p^2r^{2p-2}}{(1+r^p)^2},
\]
and each tangential Hessian eigenvalue is
\(p r^{p-2}/(1+r^p)\).
Splitting into \(r\le1\) and \(r\ge1\) proves the bounds; the formulas
extend continuously at zero. No dimension enters these estimates.

If \(\upsilon\ge1/2\), every child was shrunk. The denominator
comparison \eqref{pred:denominator-lower-bound}, the sum of
\eqref{pred:shrunk-numerator-motion}, and
\eqref{pred:log-derivatives} yield
\begin{equation}
 \log\beta_J^0(z)-\log\widetilde\beta(b)
       \le CpK_*\varepsilon\upsilon,\qquad \upsilon\ge1/2.
 \label{pred:balanced-comparison}
\end{equation}
The factor \(\upsilon\) here costs only an absolute factor of two.

For \(0<\upsilon<1/2\), retaining that factor requires the unshrunk
path vector. Write \(\sigma_a=\lambda_a/\abs{\lambda_a}\) and remove
\(\mathcal D_a(b)\) and \(\sigma_a\) from the predicted numerator.
With
\begin{align*}
 Y&=l_a^{1/p}X_a,\\
 C_a&=l_a^{1-\mu-1/p}\abs{\lambda_a}
      =l_a^{1-\alpha/p}(n_a/m_a)^\mu,\\
 Z_1(b)&=\sigma_a^{-1}\mathcal D_a(b)^{-1}
        \sum_{c\ne a}l_c^{1-\mu}\lambda_c
                         \mathcal D_c(b)\widetilde X_c,\\
 R_1&=\sum_{c\ne a}l_c\norm{\widetilde X_c}^p,
\end{align*}
the prediction becomes
\begin{equation}
 \widetilde\beta(b)=
       \frac{1+\norm{C_aY+Z_1(b)}^p}
            {1+\norm Y^p+R_1}.
 \label{pred:path-baseline}
\end{equation}
Here \(Y,C_a\), and the denominator are constant on the whole line.
Since \(n_a\le m_a\le n\), both \(l_a\) and \(n_a/m_a\) lie in
\([1-\upsilon,1]\). As \(1-\alpha/p\ge0\), we have
\begin{equation}
 0<C_a\le1,\qquad 1-C_a\le C\upsilon,\qquad
 \norm{Z_1(b)}\le K_*\upsilon.
 \label{pred:path-coefficient}
\end{equation}

In the actual numerator remove \(\mathcal D_a(b)\) and
\(\sigma_a'=\lambda_a'/\abs{\lambda_a'}\). Define
\[
 Y'=l_a'^{1/p}\mathcal D_a(b)^{-1}X_a',\qquad
 C_a'=l_a'^{1-\mu-1/p}\abs{\lambda_a'},\qquad
 R_1'=\sum_{c\ne a}l_c'\norm{X_c'}^p,
\]
and let \(Z_1'\) be the remaining, similarly rotated off-path sum.
The actual ratio has the form \eqref{pred:path-baseline} with primes.
We have
\begin{align}
 \norm{Y'-Y}&\le C\varepsilon K_*,
 &\abs{C_a'-C_a}&\le C\varepsilon\upsilon,\notag\\
 \norm{Z_1'-Z_1}&\le C\varepsilon K_*\upsilon,
 &R_1'&\ge R_1.
 \label{pred:path-errors}
\end{align}
The first and last statements follow from
\eqref{pred:scaled-motion} and \eqref{pred:denominator-lower-bound}.
For the off-path sum, sum \eqref{pred:shrunk-numerator-motion} over
\(c\ne a\), whose weights sum to at most \(\upsilon\).
Changing the removed scalar phase costs at most \(C\varepsilon\)
times the off-path size, because
\(\abs{\sigma_a'-\sigma_a}\le C\varepsilon\).

The extra \(\upsilon\) in the coefficient estimate follows from the
complementary-mass assertion of \Cref{hier:mass-motion}.
If \(r=A/(A+B)\), with \(A>0\), \(B\ge0\), and the positive
summands change relatively by at most \(C\varepsilon\), then
\[
 |r'-r|\le C\varepsilon r(1-r),\qquad
 |\log r'-\log r|\le C\varepsilon(1-r).
\]
If \(B=0\), the ratio is identically one.
Apply this to \(n_a/n\), with complement \(n-n_a\), and to
\(n_a/m_a\), with complement \(m_a-n_a\).
Both values of \(1-r\) are at most \(\upsilon\).
The bounded nonnegative powers defining \(C_a\) give its claimed
variation. This concerns the modulus of the path weight; its scalar
phase has already been removed.

Set \(\mathfrak r=R_1/(1+\norm Y^p+R_1)\). We claim
\begin{align}
 \log\beta_J^0(z)-\log\widetilde\beta(b)
    &\le C p^2K_*^2\varepsilon\upsilon
           +C pK_*\varepsilon\mathfrak r,\notag\\
 -\log\widetilde\beta(b)
    &\ge \mathfrak r-CpK_*\upsilon.
 \label{eq:src-26}
\end{align}
Hold \(Z_1=Z_1(b)\) fixed when comparing \(Y\) with \(Y'\).
The derivative in \(V\) of
\(\mathcal L(C_aV+Z_1)-\mathcal L(V)\), along that segment, is bounded by
\[
 p|1-C_a|+Cp^2\bigl(|1-C_a|\norm V+\norm{Z_1}\bigr)
       \le Cp^2K_*\upsilon.
\]
Equation~\eqref{pred:path-errors} bounds its change by
\(Cp^2K_*^2\varepsilon\upsilon\).
The changes in \(C_a,Z_1\) fit the same bound by
\eqref{pred:log-derivatives}.

For the residual denominator, set
\[
 J(V)=-\log(1+R_1e^{-\mathcal L(V)}).
\]
Its derivative has norm at most
\(pR_1/(e^{\mathcal L(V)}+R_1)\).
On the segment from \(Y\) to \(Y'\),
\(\abs{\mathcal L(V)-\mathcal L(Y)}\le CpK_*\varepsilon\).
This tends to zero as \(L\) increases, so that residual ratio is at
most \(2\mathfrak r\), and
\(J(Y')-J(Y)\le CpK_*\varepsilon\mathfrak r\).
The actual residual \(R_1'\ge R_1\) can only decrease the actual
ratio. This proves the first inequality of \eqref{eq:src-26}.
For the second, use \(C_a\le1\), \eqref{pred:log-derivatives},
and \(\log(1+t)\ge t/(1+t)\):
\begin{align*}
 -\log\widetilde\beta
 &=\mathcal L(Y)-\mathcal L(C_aY+Z_1)
       +\log(1+R_1e^{-\mathcal L(Y)})\\
 &\ge-p\norm{Z_1}+\mathfrak r
 \ge\mathfrak r-CpK_*\upsilon.
\end{align*}

\subsection{Absorbing the baseline error by blending}

The parameter sizes give
\begin{equation}
 pK_* = G^{o(1)}\quad\text{as }L\longrightarrow\infty,
 \label{pred:subpower-amplitude}
\end{equation}
with every preceding constant fixed.
Indeed \(\log N\le2L\), \(\mu=O(L^{-1/2})\), and
\(\log G\asymp L^{7/10}\), so
\(\log(pK_*)=O(\sqrt L+\log L)=o(\log G)\).
All the preceding smallness requirements consequently follow from a
single sufficiently large lower bound on \(L\).

Set \(\theta_{\mathrm{bl}}=\sqrt\varepsilon\) and blend the prediction with one:
\begin{equation}
 \mathcal B(b)=(1-\theta_{\mathrm{bl}})\widetilde\beta(b)+\theta_{\mathrm{bl}}.
 \label{pred:blend}
\end{equation}
It retains average at most one for every probability in
\Cref{pred:fourier-average}. Concavity of the logarithm gives
\[
 \log \mathcal B-\log\widetilde\beta
       \ge-\theta_{\mathrm{bl}}\log\widetilde\beta.
\]
For \(\upsilon<1/2\), Equation~\eqref{eq:src-26} yields
\begin{align*}
 \log\beta_J^0-\log \mathcal B
 &\le\bigl(Cp^2K_*^2\varepsilon+CpK_*\theta_{\mathrm{bl}}\bigr)\upsilon
       +\bigl(CpK_*\varepsilon-\theta_{\mathrm{bl}}\bigr)\mathfrak r\\
 &\le\bigl(Cp^2K_*^2\varepsilon
                  +CpK_*\sqrt\varepsilon\bigr)\upsilon.
\end{align*}
The last coefficient of \(\mathfrak r\) is nonpositive for large \(L\).
For \(\upsilon\ge1/2\), use \eqref{pred:balanced-comparison} and
\(\log\widetilde\beta\le\mathcal L(F)\le p\norm F\le CpK_*\)
to obtain the same last bound. Therefore
\begin{equation}
 \beta_J^0(z)\le e^{a_L\upsilon}\mathcal B(b),\qquad
 a_L=Cp^2K_*^2\varepsilon+CpK_*\sqrt\varepsilon
       \le G^{-\gamma/3}
 \label{pred:blended-domination}
\end{equation}
for sufficiently large \(L\).
The baseline has absorbed the term proportional to \(\mathfrak r\);
the remaining error is proportional to \(\upsilon\).

\subsection{Changing to uniform measure without a loss per level}

Choose a fixed nonnegative Schwartz function \(k\) of integral one
with \(\supp\widehat k\subset[-1/4,1/4]\).
For example, square the absolute value of the inverse Fourier
transform of a nonzero smooth function supported in \((-1/8,1/8)\)
and normalize its integral.
Let \(H=H_j\), \(h=G^{-1/16}\), \(s=HG^{-1/8}\), and
\(k_s(t)=s^{-1}k(t/s)\).
Let \(I^+\) be the concentric enlargement of \(I\) of length
\((1+h)H\), and define
\begin{equation}
 \dd\nu_I(t)=
      \left(\frac{\one_{I^+}}{|I^+|}*k_s\right)(t)\,\dd t.
 \label{pred:majorant-probability}
\end{equation}
This is a positive probability whose Fourier transform is supported
in \([-1/(4s),1/(4s)]\).
Since \(H\delta=\epsilon'G^{1/4}\), this support lies strictly inside
\((-\delta,\delta)\) for large \(L\).

Every point of \(I\) is at distance at least \(hH/2\) from the
complement of \(I^+\). The Schwartz tail beyond
\((hH)/(2s)=G^{1/16}/2\) is \(O(G^{-1/8})\), by a fixed tail estimate.
Thus the density in \eqref{pred:majorant-probability} is at least
\(H^{-1}(1-CG^{-1/16})\) on \(I\).
Writing \(\omega_I\) for uniform probability on \(I\), choose
\(0<\eta\le CG^{-1/16}<1/2\) so that
\begin{equation}
 \nu_I\ge(1-\eta)\omega_I,\qquad
 \nu_I=(1-\eta)\omega_I+\eta\nu_{\rm rem}
 \label{pred:measure-decomposition}
\end{equation}
for a positive probability \(\nu_{\rm rem}\).

The bound \(\int \mathcal B\,\dd\nu_I\le1\) alone would lose an amount of
order \(\eta\) on changing measure. We need a bound proportional to
\(\upsilon\).
If \(\widetilde\beta\le1\) on the whole line, then \(\mathcal B\le1\) and
its uniform average is already at most one.
Otherwise choose \(b_0\) with \(\widetilde\beta(b_0)>1\).
When \(\upsilon<1/2\), the denominator and path vector in
\eqref{pred:path-baseline} are fixed. For all \(b\in\R\),
\begin{align*}
 \log\widetilde\beta(b)
 &\ge\log\widetilde\beta(b_0)
          -p\norm{Z_1(b)-Z_1(b_0)}\\
 &\ge-2pK_*\upsilon.
\end{align*}
Consequently
\begin{equation}
 (1-\widetilde\beta(b))_+\le CpK_*\upsilon
       \qquad(b\in\R).
 \label{pred:negative-deviation}
\end{equation}
This does not assume an upper bound on \(R_1\): its contribution
cancels in the logarithmic difference. Equivalently, a value exceeding
one already forces
\(R_1/(1+\norm Y^p)\le e^{CpK_*\upsilon}-1\).
For \(\upsilon\ge1/2\), Equation~\eqref{pred:negative-deviation}
follows from nonnegativity since \(pK_*\upsilon\ge1\).
The blend inherits the same lower bound, as
\((1-\mathcal B)_+=(1-\theta_{\mathrm{bl}})(1-\widetilde\beta)_+\).

In the case where the prediction exceeds one somewhere,
\eqref{pred:measure-decomposition} gives
\[
 1\ge\int \mathcal B\,\dd\nu_I
 \ge(1-\eta)\int \mathcal B\,\dd\omega_I
           +\eta(1-CpK_*\upsilon).
\]
Including the easier case \(\mathcal B\le1\), this proves
\begin{equation}
 \int \mathcal B\,\dd\omega_I\le1+C G^{-1/16}pK_*\upsilon.
 \label{pred:uniform-average}
\end{equation}
The relative mass drop has survived the change of probability.

\subsection{Step supremums and normalization}

The points \(s_R,t_c\) lie in \(J\), so all component frequencies
of \(\mathcal D_c\) have magnitude at most \(C\rho_j/w\).
For \(\upsilon\ge1/2\), differentiating \eqref{pred:model} and using
\eqref{pred:snapshot-size} gives
\(\norm{\partial_b F(b)}\le C(\rho_j/w)K_*\).
For \(\upsilon<1/2\), remove the path carrier as in
\eqref{pred:path-baseline}. The off-path frequencies are
\((t_a-t_c)\eta_\xi\), again of magnitude at most \(C\rho_j/w\), so
\(\norm{\partial_b Z_1(b)}\le C(\rho_j/w)K_*\upsilon\).
In both cases,
\[
 \left|\frac{\dd}{\dd b}\log\widetilde\beta(b)\right|
       \le CpK_*\upsilon\frac{\rho_j}{w}.
\]
Blending cannot increase the absolute logarithmic derivative, because
\[
 \frac{\dd}{\dd b}\log \mathcal B
   =\frac{(1-\theta_{\mathrm{bl}})\widetilde\beta}{\mathcal B}
               \frac{\dd}{\dd b}\log\widetilde\beta,\qquad
 0\le\frac{(1-\theta_{\mathrm{bl}})\widetilde\beta}{\mathcal B}\le1.
\]

On each \(H_{j-1}\)-subcell \(I'\subset I\), let \(\mathcal B^\sharp\) equal
\(\sup_{b\in\overline{I'}}\mathcal B(b)\).
The prediction is positive and continuous, and these compact-cell
supremums are finite. The meshes are nested. Since
\(H_{j-1}\rho_j/w=G^{-1/2}\),
\begin{equation}
 \mathcal B\le \mathcal B^\sharp\le e^{b_L\upsilon}\mathcal B,\qquad
 b_L=CpK_*G^{-1/2}.
 \label{pred:step-supremum}
\end{equation}
Use the fixed half-open mesh convention at endpoints; taking
supremums over closed subcells preserves the inequalities.
Define
\begin{equation}
 m_I=\max\left(1,\int \mathcal B^\sharp\,\dd\omega_I\right),\qquad
 \Gamma_j=\mathcal B^\sharp/m_I\quad\text{on }I.
 \label{pred:gamma-normalization}
\end{equation}
This function is nonnegative, constant on the finer cells, and has
average at most one. Equations~\eqref{pred:uniform-average} and
\eqref{pred:step-supremum} yield
\begin{equation}
 \log m_I\le\bigl(b_L+CG^{-1/16}pK_*\bigr)\upsilon.
 \label{pred:normalization-cost}
\end{equation}

\begin{proof}[Proof of \Cref{pred:marked-comparison}]
Perform the construction on each marked cell meeting \(Z\).
Equations~\eqref{pred:blended-domination},
\eqref{pred:step-supremum}, and \eqref{pred:gamma-normalization} give
\[
 \beta_J^0(z)\le e^{a_L\upsilon}\mathcal B(b)
       \le e^{a_L\upsilon}\mathcal B^\sharp(b)
       =e^{a_L\upsilon}m_I\Gamma_j(b).
\]
By \eqref{pred:normalization-cost}, the exponent is at most
\[
 \bigl(G^{-\gamma/3}+CpK_*G^{-1/2}
                    +CpK_*G^{-1/16}\bigr)\upsilon.
\]
Equations~\eqref{pred:subpower-amplitude} and
\eqref{pred:drop-comparability} bound this by
\(G^{-c_{\rm av}}u_j(z)\), for a fixed positive exponent:
choose \(c_{\rm av}<\min(\gamma/4,1/64)\) and then take \(L\)
sufficiently large. This proves \eqref{eq:src-24} with no error
independent of the mass drop.

There are countably many cells on each fixed line and finitely many
levels in the current truncation. Choose one snapshot in every
relevant cell, use \eqref{pred:gamma-normalization} there, and set
\(\Gamma_j=1\) elsewhere. This is a measurable step function of the
line parameter. The proof of \Cref{res:path-product} integrates the
original integrand on lines where its slice is measurable, as holds for
almost every transverse parameter. It therefore needs no jointly
measurable snapshot choices across parallel lines.
\end{proof}
